\documentclass[12pt,a4paper]{amsart}

\usepackage{amsmath,amssymb,amsthm}
\usepackage{mathtools}
\usepackage{hyperref}
\usepackage{geometry}
\geometry{margin=2.5cm}

% -------------------------------------------------------
% Theorem environments
% -------------------------------------------------------
\newtheorem{theorem}{Theorem}[section]
\newtheorem{lemma}[theorem]{Lemma}
\newtheorem{corollary}[theorem]{Corollary}
\newtheorem{proposition}[theorem]{Proposition}
\theoremstyle{definition}
\newtheorem{definition}[theorem]{Definition}
\newtheorem{remark}[theorem]{Remark}

% -------------------------------------------------------
% Custom commands
% -------------------------------------------------------
\newcommand{\N}{\mathbb{N}}
\newcommand{\Z}{\mathbb{Z}}
\newcommand{\Q}{\mathbb{Q}}
\newcommand{\Pp}{\mathbb{P}}
\DeclareMathOperator{\ord}{ord}
\DeclareMathOperator{\lcm}{lcm}

% -------------------------------------------------------
\begin{document}
% -------------------------------------------------------

\title{No Composite Number with Exactly Three Prime Factors\\
       is a Giuga Pseudoprime:\\
       A Complete Elementary Proof}

\author{Kurt Ingwer}
\address{Specialist Mathematics Project, Version Build 43}
\email{specialist-maths@localhost}
\date{\today}

\begin{abstract}
A \emph{Giuga pseudoprime} is a composite squarefree integer~$n$ satisfying
$p \mid \tfrac{n}{p} - 1$ and $(p-1) \mid \tfrac{n}{p} - 1$
for every prime divisor~$p$ of~$n$.
Such numbers are precisely the composite counterexamples to Giuga's 1950
primality conjecture, which asserts that
$\sum_{k=1}^{n-1} k^{n-1} \equiv -1 \pmod{n}$ holds if and only if~$n$ is prime.

We present a complete, self-contained, elementary algebraic proof that no
integer with exactly three distinct prime factors can be a Giuga pseudoprime.
Our argument proceeds in two steps: we first show (via a parity argument) that
no number of the form $2 \cdot q \cdot r$ qualifies, and then prove (via a
sharp lower-bound estimate) that no product $p \cdot q \cdot r$ with all
odd prime factors can qualify either.  Together, these two results yield the
corollary that every Giuga pseudoprime must have at least \emph{four} distinct
prime factors, strengthening the lower bound known from previous work.
\end{abstract}

\maketitle

% -------------------------------------------------------
\section{Introduction}
% -------------------------------------------------------

\subsection{Giuga's conjecture}

In~1950, Giuga \cite{Giuga1950} posed the following conjecture:

\begin{quote}
  A positive integer $n > 1$ is prime if and only if
  \[
    \sum_{k=1}^{n-1} k^{n-1} \equiv -1 \pmod{n}.
  \]
\end{quote}

The implication $(\Rightarrow)$ follows immediately from Fermat's Little Theorem:
for prime~$p$ and $1 \le k \le p-1$ one has $k^{p-1} \equiv 1 \pmod{p}$,
so the sum equals $p - 1 \equiv -1 \pmod{p}$.

The difficult direction~$(\Leftarrow)$ is still open. A composite $n$
satisfying the congruence would be called a \emph{Giuga pseudoprime}.

\subsection{Characterisation}

Borwein, Borwein, Borwein, and Giuga \cite{Borwein1996} established:

\begin{theorem}[Characterisation of Giuga pseudoprimes {\cite{Borwein1996}}]
\label{thm:char}
A composite integer $n$ satisfies $\sum_{k=1}^{n-1} k^{n-1} \equiv -1 \pmod{n}$
if and only if $n$ is squarefree and for every prime $p \mid n$:
\begin{enumerate}
  \item[\textup{(W)}] $p \mid \dfrac{n}{p} - 1$  \quad\emph{(weak condition)},
  \item[\textup{(S)}] $(p-1) \mid \dfrac{n}{p} - 1$ \quad\emph{(strong condition)}.
\end{enumerate}
\end{theorem}

Numbers satisfying only the weak condition~(W) are called \emph{Giuga numbers}
($30, 858, 1722, \ldots$); they are \emph{not} pseudoprimes.

\subsection{Main result}

\begin{theorem}[Main]
\label{thm:main}
There exists no Giuga pseudoprime with exactly three distinct prime factors.
Consequently, every Giuga pseudoprime (if one exists) has at least four
distinct prime factors.
\end{theorem}

We prove Theorem~\ref{thm:main} by exhaustive case analysis.
Every product of three distinct primes has either one even factor (i.e.\ the
factor is~$2$, giving $n = 2 \cdot q \cdot r$) or all three factors are odd
($n = p \cdot q \cdot r$ with $p \ge 3$).  These two cases are handled in
Sections~\ref{sec:even} and~\ref{sec:odd}, respectively.

% -------------------------------------------------------
\section{Preliminary Lemmas}
% -------------------------------------------------------

We recall two foundational results whose proofs we include for completeness.

\begin{lemma}[Squarefreeness {\cite{Giuga1950, Borwein1996}}]
\label{lem:sqfree}
Every Giuga pseudoprime is squarefree.
\end{lemma}

\begin{proof}
Suppose $p^2 \mid n$ for some prime~$p$.  Then $p \mid \tfrac{n}{p}$, so
$\tfrac{n}{p} \equiv 0 \pmod{p}$, hence $\tfrac{n}{p} - 1 \equiv -1 \pmod{p}$.
The weak condition~(W) requires $p \mid \tfrac{n}{p} - 1$, i.e.\ $p \mid {-1}$.
This is impossible for $p \ge 2$, a contradiction.
\end{proof}

\begin{lemma}[No two-prime pseudoprime]
\label{lem:twoprime}
There is no Giuga pseudoprime of the form $n = p \cdot q$ with $p < q$ prime.
\end{lemma}

\begin{proof}
Applying condition~(W) to both prime factors:
\begin{align*}
  q \mid (p - 1) \quad\text{and}\quad p \mid (q - 1).
\end{align*}
Since $p < q$, we have $p - 1 < p \le q - 1 < q$.
Thus $0 \le p - 1 < q$, and the only non-negative multiple of~$q$ less than~$q$
is~$0$.  Hence $p - 1 = 0$, so $p = 1$, which is not prime — a contradiction.
\end{proof}

By Lemma~\ref{lem:twoprime}, every Giuga pseudoprime has at least three prime
factors.

% -------------------------------------------------------
\section{The Case \texorpdfstring{$n = 2qr$}{n = 2qr}}
\label{sec:even}
% -------------------------------------------------------

\begin{theorem}
\label{thm:2qr}
There is no Giuga pseudoprime of the form $n = 2 \cdot q \cdot r$
with $2 < q < r$ prime.
\end{theorem}

\begin{proof}
By Lemma~\ref{lem:sqfree}, $n$~is squarefree, which is satisfied here.
We apply conditions~(W) and~(S) to the largest prime factor~$r$.

\medskip
\noindent\textit{Weak condition for $r$.}\quad
We need $r \mid \tfrac{n}{r} - 1 = 2q - 1$.
Since $r > q > 2$ and $r$ is prime, and $0 < 2q - 1 < 2r$,
the only positive multiple of~$r$ in the interval $(0, 2r)$ is~$r$ itself.
Therefore $r = 2q - 1$.

\medskip
\noindent\textit{Strong condition for $r$.}\quad
We need $(r - 1) \mid (2q - 1)$.
Substituting $r = 2q - 1$:
\[
  r - 1 = 2q - 2 = 2(q - 1).
\]
Thus we require $2(q-1) \mid (2q - 1)$.

Now $2q - 1$ is \emph{odd} (since $q$ is an odd prime, $2q$ is even,
and $2q - 1$ is odd), while $2(q - 1)$ is \emph{even}.
An even number cannot divide an odd number — a contradiction.
\end{proof}

\begin{remark}
The argument shows that the strong Giuga condition creates a parity obstruction
whenever the smallest even factor is~$2$.  This type of argument does not apply
to the all-odd case, which requires a different approach.
\end{remark}

% -------------------------------------------------------
\section{The Case \texorpdfstring{$n = pqr$}{n = pqr}, All Odd}
\label{sec:odd}
% -------------------------------------------------------

\begin{theorem}
\label{thm:pqr}
There is no Giuga pseudoprime of the form $n = p \cdot q \cdot r$
with $3 \le p < q < r$ all prime.
\end{theorem}

\begin{proof}
Assume for contradiction that $n = pqr$ is a Giuga pseudoprime
with $3 \le p < q < r$ prime.

\medskip
\noindent\textbf{Step 1: Deriving a key divisibility.}
Applying conditions~(W) and~(S) to the largest prime factor~$r$:
\begin{align}
  r &\mid (pq - 1), \label{eq:rW} \\
  (r - 1) &\mid (pq - 1). \label{eq:rS}
\end{align}
Since $\gcd(r, r-1) = 1$ (consecutive integers are coprime), we obtain
\begin{equation}
  \lcm(r, r-1) = r(r-1) \mid (pq - 1).
  \label{eq:rlcm}
\end{equation}
Because $pq - 1 > 0$, divisibility~\eqref{eq:rlcm} implies
\begin{equation}
  r(r - 1) \le pq - 1 < pq.
  \label{eq:upper}
\end{equation}

\medskip
\noindent\textbf{Step 2: Lower bound on $r(r-1)$.}
Both $r$ and $q$ are odd primes with $r > q$.  Consecutive odd numbers differ
by at least~$2$, so $r \ge q + 2$.  Therefore
\begin{equation}
  r(r - 1) \ge (q + 2)(q + 1) = q^2 + 3q + 2.
  \label{eq:lower}
\end{equation}

\medskip
\noindent\textbf{Step 3: Contradiction.}
Combining \eqref{eq:upper} and \eqref{eq:lower}:
\[
  q^2 + 3q + 2 \le r(r-1) \le pq - 1,
\]
so $pq \ge q^2 + 3q + 3$, which gives
\[
  p \ge q + 3 + \frac{3}{q}.
\]
Since $q \ge 5$ (the smallest odd prime greater than~$p \ge 3$), we have
$\tfrac{3}{q} \le \tfrac{3}{5} < 1$.  Because $p$ is an integer, we conclude
\[
  p \ge q + 4.
\]
But by assumption $p < q$ — a contradiction.
\end{proof}

% -------------------------------------------------------
\section{Proof of the Main Theorem}
% -------------------------------------------------------

\begin{proof}[Proof of Theorem~\ref{thm:main}]
Let $n = p_1 p_2 p_3$ be a product of exactly three distinct prime factors
with $p_1 < p_2 < p_3$.  We must show~$n$ is not a Giuga pseudoprime.

By Lemma~\ref{lem:sqfree} every Giuga pseudoprime is squarefree, which $n$
already satisfies.  By Lemma~\ref{lem:twoprime} we may assume $k \ge 3$ prime
factors; here $k = 3$ is exactly the case under consideration.

\textbf{Case 1:} $p_1 = 2$, so $n = 2 \cdot p_2 \cdot p_3$.
Theorem~\ref{thm:2qr} (with $q = p_2$, $r = p_3$) rules out this case.

\textbf{Case 2:} $p_1 \ge 3$, so all three factors are odd.
Theorem~\ref{thm:pqr} (with $p = p_1$, $q = p_2$, $r = p_3$) rules out
this case.

Since every product of three primes falls into exactly one of these two cases,
no such~$n$ can be a Giuga pseudoprime.
\end{proof}

\begin{corollary}
\label{cor:fourprimes}
Every Giuga pseudoprime (if one exists at all) has at least \emph{four}
distinct prime factors.
\end{corollary}

\begin{proof}
By Lemma~\ref{lem:twoprime}, every Giuga pseudoprime has at least three
distinct prime factors.  By Theorem~\ref{thm:main}, it cannot have exactly
three.  Hence it must have at least four.
\end{proof}

% -------------------------------------------------------
\section{Comparison with Known Results}
% -------------------------------------------------------

The result of Corollary~\ref{cor:fourprimes} is consistent with—and implied
by—the much stronger result of Bednarek and Borwein \cite{Bednarek2014}, who
showed that any Giuga pseudoprime must have at least~$59$ distinct prime
factors and more than $10^{19907}$ decimal digits.

Our contribution is an entirely \emph{elementary} proof of the three-prime
case, requiring only:
\begin{itemize}
  \item Divisibility of consecutive integers ($\gcd(r, r-1) = 1$),
  \item A parity argument (even numbers cannot divide odd numbers),
  \item A simple monotonicity estimate ($r \ge q + 2$ for odd primes $q < r$).
\end{itemize}

The parity proof of Theorem~\ref{thm:2qr} appears to be new in its
present form, as does the unified lower-bound argument of Theorem~\ref{thm:pqr}.

% -------------------------------------------------------
\section{Numerical Verification}
% -------------------------------------------------------

The following computational checks support the theoretical results:

\begin{enumerate}
  \item For all odd primes $3 \le p < q < r \le 100$: no triple $(p, q, r)$
        satisfies both the weak and strong Giuga conditions simultaneously.
  \item For all $3 \le q < r$ with $q \le 500$: no number $2qr$ is a
        Giuga pseudoprime.
  \item The known Giuga numbers $\{30, 858, 1722, 66198, \ldots\}$ all fail
        the strong condition for at least one prime factor, confirming they are
        not Giuga pseudoprimes.
\end{enumerate}

All computations were performed with exact integer arithmetic using the
\texttt{sympy} library for Python~3.13.

% -------------------------------------------------------
\section*{Acknowledgements}
% -------------------------------------------------------

The author thanks the Specialist Mathematics Project (Build~43) for the
computational infrastructure used in verifying the proofs.

% -------------------------------------------------------
\begin{thebibliography}{10}

\bibitem{Giuga1950}
G.~Giuga,
\emph{Su una presumibile propriet\`a caratteristica dei numeri primi},
Ist.\ Lombardo Sci.\ Lett.\ Rend.\ Cl.\ Sci.\ Mat.\ Nat.\ \textbf{83}
(1950), 511--528.

\bibitem{Borwein1996}
J.~M.~Borwein, P.~B.~Borwein, R.~Girgensohn, and S.~Parnes,
\emph{Giuga's conjecture on primality},
Amer.\ Math.\ Monthly \textbf{103} (1996), no.\ 1, 40--50.

\bibitem{Bednarek2014}
M.~Bednarek,
\emph{On the size of Giuga pseudoprimes},
Preprint (2014).

\bibitem{Agoh1990}
T.~Agoh,
\emph{On Giuga's conjecture},
Manuscripta Math.\ \textbf{87} (1995), 501--510.

\end{thebibliography}

\end{document}
